\subsection{The original Hermitian class sum from convergent Siegel--Weil}
\label{actual-hermitian-siegel-weil}

Let $E/F$ be the CM extension of an original Host. In particular,
$E/F$ is unramified at every finite place. Retain its prescribed CM
type $\Psi$, the positive Hermitian space $V=(E^n,1)$, and the
source's $\mathrm{Vset}(E,n,M)$, $\mathrm{DH}(E,n,M)$ and
$\mathrm{classSum}$. Thus a lattice is self-dual for the
$\mathcal O_E$-valued Hermitian form, is standard at primes
dividing the positive integer $M$, and has weight
\[
 w_M(\Lambda)=|\operatorname{Aut}_M(\Lambda)|^{-1},\qquad
 \operatorname{Aut}_M(\Lambda)=
 \{h\in\mathrm{DH}(E,n,M):h\Lambda=\Lambda\}.
\]
These are the original weights, without a mass normalization.

\begin{theorem}\label{actual-hermitian-siegel-weil-theorem}
For every even $n$, every $N$ with $4N<n$, all positive integers
$M,R$, and every original admitted family
$\varphi=(\varphi_\Lambda)_{\Lambda\in\mathrm{Vset}(E,n,M)}$
of degree $N$, one has, for every $Z\in\mathfrak H_N^\Psi$,
\begin{equation}\label{actual-sw-class-identity}
 \sum_{[\Lambda]}w_M(\Lambda)
 \bigl(\theta_{\varphi_\Lambda}(Z,0,0)
                   -E_{\varphi_\Lambda}(Z,0,0)\bigr)=0.
\end{equation}
The theta series, Eisenstein terms, rational Weil action and
indexing of rational maximal isotropic spaces are literally those
in \texttt{FirstCoordinateAdmissible}. All sums converge absolutely.
This supplies the original \texttt{SWAt} instance for every
actual Host, with all its $M,R,\varphi,Z$ quantifiers.
\end{theorem}

\begin{proof}
If $N=0$, both series have a single term, so assume $N>0$.
Put $H=U(V)$ and let $K_M\subset H(\mathbb A_{F,f})$ be the
standard lattice stabilizer, with principal congruence condition
$h\equiv1\pmod M$ at primes dividing $M$.

\textbf{The exact lattice classes.}
Every self-dual lattice in $V_v$ is isometric to
$\mathcal O_{E,v}^n$. At a split place the two summands are
mutually dual, so choose a basis in one and the dual basis in the
other. At an unramified field place, a nondegenerate residual
Hermitian form has a vector of nonzero norm. Indeed, if all
diagonal norms vanished, polarization with a scalar of nonzero
residual trace would contradict nondegeneracy. Lift that vector,
split off its unit-norm line, and continue by orthogonal
induction. Every resulting diagonal unit is a norm of a unit in
the unramified quadratic extension. Rescaling gives the standard
lattice. This proof includes residue characteristic two
\cite{actual-sw-unit-norms}.

The group $H$ has weak approximation at any finite set of places.
For completeness, its Cayley chart
\[
             X^*=-X,\qquad c(X)=(1-X)(1+X)^{-1}
\]
is a rational open chart containing $1$. The excluded determinant
locus has empty interior over every local field; approximate the
given local elements inside the chart and then approximate their
skew-Hermitian coordinates over $F$. No integrality away from the
specified finite set is asserted or needed.

For every $\Lambda\in\mathrm{Vset}$ choose $h_\Lambda$ carrying
$\widehat{\mathcal O}_E^n$ to $\widehat\Lambda$, and equal to $1$
at the primes of $M$. Conversely, weak approximation makes any
adelic representative principal congruent to $1$ at those primes
by a left rational unitary translation. Its lattice is then in
the original $\mathrm{Vset}$. Two representatives give the same
double coset exactly when their lattices differ by a rational
unitary matrix congruent to $1$ modulo $M$ at those primes.
Clearing its finitely many remaining denominators by an integer
prime to $M$ identifies that condition with the original
definition of $\mathrm{DH}$. Consequently
\begin{equation}\label{actual-sw-class-bijection}
 \mathrm{Cls}(E,n,M)\simeq
              H(F)\backslash H(\mathbb A_{F,f})/K_M,
 \qquad
 H(F)\cap h_\Lambda K_Mh_\Lambda^{-1}
                  =\operatorname{Aut}_M(\Lambda).
\end{equation}
The class set is finite and each stabilizer is finite. One
elementary verification uses the positive integral trace form on
$\Lambda$: its rank and determinant are fixed, the latter by
$|\operatorname{disc}E|^n$. Minkowski reduction gives finitely many
integral Gram matrices. For each such matrix, multiplication by
fixed generators of $\mathcal O_E$ has bounded operator norm,
hence only finitely many possible integral matrices. This gives
finitely many Hermitian lattice classes. Their possible
$M$-markings form a finite set. An integral automorphism of a
positive form lies in a compact group and hence is one of
finitely many integral matrices.

\textbf{An admitted family is one finite Schwartz vector.}
The original support and period conditions are
\[
 \operatorname{supp}\varphi_\Lambda\subset R^{-1}\Lambda^N,
 \qquad \varphi_\Lambda(x+Ry)=\varphi_\Lambda(x)
                         \quad(y\in\Lambda^N).
\]
They extend $\varphi_\Lambda$ uniquely to a locally constant,
compactly supported function $\varphi_{\Lambda,f}$ on
$V(\mathbb A_{F,f})^N$. Write $\varphi_f$ for the standard
lattice's extension. The full admitted contract gives
\begin{equation}\label{actual-sw-family-equiv}
 \varphi_{\Lambda,f}(x)=\varphi_f(h_\Lambda^{-1}x),
 \qquad \varphi_f(k^{-1}x)=\varphi_f(x)\quad(k\in K_M).
\end{equation}
Here is the locality argument, including its global integrality
point. Approximate $h_\Lambda$ at the primes of $RM$ by
$q\in\mathrm{DH}$, closely enough to identify the two lattices
there and their $R^{-1}/R$ coefficient quotients. Then
$q\mathcal O_E^n$ and $\Lambda$ satisfy the original
$\mathrm{Agree}(RM)$ condition. For a rational $x$ in
$R^{-1}\Lambda^N$, the Chinese remainder theorem chooses
\[
 y\in R^{-1}\Lambda^N\cap R^{-1}(q\mathcal O_E^n)^N,
               \qquad y-x\in R\Lambda^N,
\]
with $q^{-1}y$ representing the same coefficient class as
$h_\Lambda^{-1}x$. At the primes of $R$ the two lattices agree;
away from $R$ the affine condition is just membership in
$\Lambda$, so these requirements are compatible. Periodicity,
admitted locality, and admitted rational equivariance now give
\[
 \varphi_\Lambda(x)=\varphi_\Lambda(y)
 =\varphi_{q\mathcal O_E^n}(y)
 =\varphi_{\mathcal O_E^n}(q^{-1}y)
 =\varphi_f(h_\Lambda^{-1}x).
\]
Outside the common support both sides are zero. Density of
rational points gives the first identity in
\eqref{actual-sw-family-equiv}. The same argument with
$\Lambda=\mathcal O_E^n$ and $h_\Lambda=k\in K_M$ gives the
second. Thus no strong approximation or new level is used.

\textbf{The original determinant and additive characters.}
Use the additive character of $\mathbb A_F/F$ whose infinite
part is $e^{2\pi i\operatorname{Tr}_{F/\mathbb Q}}$.
Choose a unitary Hecke character $\chi$ of $E$ trivial on
$\mathbb A_F^\times$ and having infinite component
\begin{equation}\label{actual-sw-infinity-type}
 \chi_\infty(a)=\prod_{\tau\in\Psi}
          \left(\frac{\overline{\tau(a)}}{|\tau(a)|}\right)^n.
\end{equation}
Such a character exists without enlarging $E$. The compact
archimedean norm-one torus embeds as a closed subgroup of the
compact group $T(F)\backslash T(\mathbb A_F)$,
$T=\ker(\operatorname{Res}_{E/F}\mathbb G_m\to\mathbb G_m)$.
Extend its character $t\mapsto\prod_\tau t_\tau^{-n/2}$ to that
compact abelian group, and compose with $a\mapsto a/\bar a$.
The character-extension theorem applies to closed subgroups
\cite{actual-sw-character-extension};
evenness of $n$ gives integral exponents. Compactness here follows
also from finiteness of ideal classes and the finite index of
unit norms, which contain the squares of $\mathcal O_F^\times$.
The resulting character is trivial on $E^\times\mathbb A_F^\times$.
Its possible finite conductor plays no role in a level or a
coefficient-field assertion.

Use the classical Weil representation for $U(N,N)\times H$
with this $\chi$; the $H$-action is exactly
$\varphi(x)\mapsto\varphi(h^{-1}x)$.
Its archimedean Gaussian is
\[
 \Phi_\infty(x)=
       \exp\left(-2\pi\sum_{\tau\in\Psi}\sum_{b,a}
                                      |\tau(x_{ba})|^2\right).
\]
The holomorphic Gaussian transforms with the exact factor
$\prod_\tau\det(C_\tau Z_\tau+D_\tau)^{-n}$.
This can be checked on the Siegel generators. A Hermitian
unipotent multiplies by the exponential of its quadratic form.
For $\operatorname{diag}(a,\bar a^{-t})$, the scalar is
\[
 \chi_\infty(\det a)|\det a|_E^{n/2}
            =\prod_\tau\overline{\det\tau(a)}^{n},
\]
which is precisely the inverse $n$th power of its denominator.
The Fourier generator uses self-dual measure $2\,d\Re z\,d\Im z$
in each complex coordinate; completing the square gives
$\det(-iZ)^{-n}$ and the oscillator phase $i^{-nN}$, whose
product is $\det Z^{-n}$. Adding the two Jacobi linear terms
in the same calculation gives exactly the original
$\mathrm{hFac}$ and the two denominator actions in
$\mathrm{TransH}$, including its $2\pi i$ exponent.

Global Poisson summation therefore says that the finite Weil
operator at every rational $\gamma$ satisfies the literal
original $\mathrm{TransH}(\gamma,\varphi,\psi)$.
That relation determines $\psi$ uniquely among Schwartz--Bruhat
functions: hold $Z$ fixed and extract the Fourier coefficients
in the independent Jacobi variables $z,z'$. Their characters
separate all rational matrices $x$, and the Gaussian series is
absolutely convergent. Thus it is the chosen source
$\mathrm{weil}(\gamma)\varphi$, not an operator differing by an
unspecified phase. The choice of the global extension in
\eqref{actual-sw-infinity-type} cannot change this conclusion;
its values on rational determinants are forced by that infinite
component.

The original finite Haar integral gives
$\widehat{\mathcal O}_E$ volume one. The self-dual trace measure
used in Poisson summation gives it volume
$|\operatorname{disc}E|^{-1/2}$ per $E$-coordinate. These are
distinct measures. The preceding comparison retains the latter
inside the Fourier operator, and does not replace the original
inner product or any original class weight.

\textbf{The probability integral and its exact class mass.}
Set
\[
       \mathfrak m_M=\sum_{[\Lambda]}w_M(\Lambda)>0.
\]
The Gaussian is invariant under the compact group $H(F_\infty)$.
The double-coset decomposition \eqref{actual-sw-class-bijection}
and \eqref{actual-sw-family-equiv} consequently give
\begin{equation}\label{actual-sw-probability-mass}
 \int_{H(F)\backslash H(\mathbb A_F)}
           \Theta(g,h;\varphi_f\Phi_\infty)\,dh_{\rm prob}
   =\mathfrak m_M^{-1}\sum_{[\Lambda]}w_M(\Lambda)
                    \Theta(g,1;\varphi_{\Lambda,f}\Phi_\infty).
\end{equation}
Indeed every component has volume the same constant divided by
$|\operatorname{Aut}_M(\Lambda)|$; total volume one fixes that
constant to $\mathfrak m_M^{-1}$.

Apply Gan--Qiu--Takeda, Theorem~7.1, with Hermitian dimension
$m=n$, split skew-Hermitian dimension $2N$, and Witt index zero.
The group $H$ is anisotropic because $V$ is positive at every
archimedean place. Their $d(N)=N$ in the unitary case, and
$n>4N$ lies in Weil's absolutely convergent range $n>2N$.
Their integral is explicitly probability normalized in
\S3.1. Their factor is
$c_{n,N}\tau(H)/[E:F]=1$: $c_{n,N}=1$ since $n>N$, and
$\tau(H)/[E:F]=1$ for a unitary group.
Thus their Eisenstein series equals the integral in
\eqref{actual-sw-probability-mass} exactly
\cite{actual-sw-gqt,actual-sw-source}.

At $g=(g_\infty,1)$ remove the same nonzero archimedean
Gaussian factor from both sides. The resulting theta series is
the original $\theta_\varphi(Z,0,0)$. The identification
$P(F)\backslash U(N,N)(F)$ with bottom-row maximal isotropic
spaces gives the original $\mathrm{Lag}$ indexing. The Gaussian
calculation identifies each summand with
$\mathrm{weil}(\gamma)\varphi(0)\mathrm{hFac}(\gamma,Z,0,0)^{-1}$,
so the resulting Eisenstein series is the original one as well.
Moreover it is independent of $\Lambda$: the finite Weil action
commutes with $H(\mathbb A_{F,f})$, and evaluation at zero is
$H$-invariant. Multiplying by $\mathfrak m_M$ now proves
\eqref{actual-sw-class-identity} with precisely its original
unnormalized class weights.
\end{proof}

Only the value at $z=z'=0$ is claimed here, as in the original
$\mathrm{SWAt}$ contract. The harmonic Jacobi comparison and
spanning argument must still use their actual heat equations,
finite test functions and reproducing scalar.
